% 注意：本片段从第 2 章（2.2.8）中间开始，父 \section 在 c03 片段中定义
% =====================================================================
%  c05：原印刷页 39–46（PDF p.40–47）
%  §2.2.8、§2.2.9，以及 §2.2.9 内一个无编号小标题
%  （Generalization To Compact Riemann Surface Of Genus g）
%  片段开头 (2.130) 及其引导句仍属上一小节 2.2.7 的收尾（承接 c04 片段末尾，
%  勿与之重复），故此处不补任何标题；本片段无图；公式编号 (2.130)–(2.180) 连续。
% =====================================================================

对于 $f(\gamma,A)=2\operatorname{Re}\chi_{fund}(\gamma,A)$、$G=GL(n,\mathbb{R})$ 或 $GL(n,\mathbb{C})$ 的情形，上式化简为

\begin{equation}
  \{f(\gamma_1,\cdot),f(\gamma_2,\cdot)\}
  = \sum_{p\in\gamma_1\cap\gamma_2}\iota_p(\gamma_1,\gamma_2)\,f\big(\gamma_1(p)\cdot\gamma_2(p),\cdot\big)
  \tag{2.130}
  \label{eq:2.130}
\end{equation}

\subsubsection{环面上平坦联络的量子化：真实版本} % 原文: 2.2.8 Quantization Of Flat Connections On The Torus: The Real Story

我们现在把讨论具体化到 $\Sigma_2=T^2$，即二维环面。

这里我们面对的是 Chern-Simons 理论量子化中最基本、也最容易上手的例子之一。因此我们会以极其详尽——乃至令人痛苦——的细致程度来处理它。

我们把它写成 $T^2=\mathbb{R}^2/\mathbb{Z}\oplus\mathbb{Z}$，也就是取坐标 $(\sigma^1,\sigma^2)$，满足

\begin{equation}
  \sigma^i\sim\sigma^i+1 \qquad i=1,2
  \tag{2.131}
  \label{eq:2.131}
\end{equation}

$\pi_1$ 有两个生成元，我们可以取下列代表：

\begin{equation}
  \begin{aligned}
    \alpha:\quad & \gamma(t)(\sigma^1(t),\sigma^2(t))=(t,0) &\quad & 0\le t\le 1\\
    \beta:\quad & \gamma(t)(\sigma^1(t),\sigma^2(t))=(0,t) &\quad & 0\le t\le 1
  \end{aligned}
  \tag{2.132}
  \label{eq:2.132}
\end{equation}

平坦 $U(1)$ 规范场的规范等价类，由绕 $\alpha$、$\beta$ 两个圈的和乐（holonomy）描述：

\begin{equation}
  U_1=\operatorname{Hol}(\alpha)=\exp i\oint_\alpha A
  \qquad\qquad
  U_2=\operatorname{Hol}(\beta)=\exp i\oint_\beta A
  \tag{2.133}
  \label{eq:2.133}
\end{equation}

我们可以取对数：

\begin{equation}
  U_i=e^{ia_i} \qquad a_i\in\mathbb{R}/2\pi\mathbb{Z}
  \tag{2.134}
  \label{eq:2.134}
\end{equation}

用微分形式写出来，$T^2$ 上最一般的平坦 $U(1)$ 规范场是

\begin{equation}
  A=a_1d\sigma^1+a_2d\sigma^2
  \tag{2.135}
  \label{eq:2.135}
\end{equation}

其中由于大规范变换，我们必须作等同 $a_i\sim a_i+2\pi\mathbb{Z}$。为了得到正确归一的辛形式，把它代入作用量：

\begin{equation}
  \frac{\kappa}{2\pi}\int_{T^2\times\mathbb{R}}A\,dA
  =\frac{\kappa}{2\pi}\int dt\,(a_1\dot a_2-a_2\dot a_1)
  \tag{2.136}
  \label{eq:2.136}
\end{equation}

于是我们重新得到上文一般结果 {\small\itshape ［原文标记：****］} 的一个特例：
% [原文待办标记，原文 p.39：「****」是作者留作公式编号占位而始终未填的记号]

\begin{equation}
  \omega=\frac{\kappa}{\pi}da_1\wedge da_2
  \tag{2.137}
  \label{eq:2.137}
\end{equation}

\textbf{评注}

\begin{enumerate}
  \item 按我们的归一化，Hilbert 空间中半经典态数由 $\omega/(2k)$ 的辛体积给出。在本情形下即
  \begin{equation}
    \int_{\mathcal{M}_{flat}}\frac{\omega}{2\pi}=2\kappa
    \tag{2.138}
    \label{eq:2.138}
  \end{equation}
  % [Tutor Note: 这里的 $2\pi$ 要核对两处。(i) (2.137) 的系数：由 (2.136) 读出辛形式时须用 $a_1\dot a_2-a_2\dot a_1=2a_1\dot a_2-\frac{d}{dt}(a_2a_1)$，于是作用量已是 $\int dt\,p\dot q$ 的形式，$p=\frac{\kappa}{\pi}a_1$、$q=a_2$，故 $\omega=dp\wedge dq=\frac{\kappa}{\pi}da_1\wedge da_2$；若漏掉这个「因子 2」而写成 $2\kappa/\pi$，(2.138) 就会差一个因子 2，(2.148) 的 $2\kappa$ 维也无从谈起。(ii) 上一行正文印作 $\omega/(2k)$（斜体 $k$），而 (2.138) 左端是 $\omega/2\pi$：按「辛体积＝半经典态数」的标准陈述此处应为 $\omega/(2\pi)$，$k$ 极可能是 $\pi$ 的误植；若把 $k$ 当作能级（即 $\kappa$），则与 $2\kappa$ 维的结论不自洽。此处照原文转写，请以 (2.138) 为准。]

  \item 要做量子化，我们必须选定相空间上哪个坐标是「位置」、哪个是「动量」。这里显然存在歧义。经典地，我们可以对 $a_i$ 作取值于 $Sp(2,\mathbb{R})$ 的辛变换。由于 $a_i$ 的周期性条件，我们只被允许作 $Sp(2,\mathbb{Z})$ 中的变换。因此取
  \begin{equation}
    \begin{pmatrix}a_1<本地路径>
    \to A\begin{pmatrix}a_1<本地路径>
    \tag{2.139}
    \label{eq:2.139}
  \end{equation}
  其中
  \begin{equation}
    A=\begin{pmatrix}\alpha & \beta\\ \gamma & \delta\end{pmatrix}
    \in SL(2,\mathbb{Z})=Sp(2,\mathbb{Z})
    \tag{2.140}
    \label{eq:2.140}
  \end{equation}

  \item 这些正则变换也可以等同于环面的大微分同胚（large diffeomorphism）的作用。\eqref{eq:2.139}、\eqref{eq:2.140} 的变换等价于
  \begin{equation}
    \begin{pmatrix}\tilde\sigma^1\\ \tilde\sigma^2\end{pmatrix}
    =\begin{pmatrix}\delta & -\gamma\\ -\beta & \alpha\end{pmatrix}
    \begin{pmatrix}\sigma^1\\ \sigma^2\end{pmatrix}
    \tag{2.141}
    \label{eq:2.141}
  \end{equation}
  显然，这是环面的一个微分同胚。事实上，按下述意义它是一个非平凡的微分同胚：

  一般地，若 $X$ 是光滑流形，可以在 $X$ 的全部微分同胚构成的群上赋予一个拓扑。这是一个无穷维（Banach）Lie 群。单位元的连通分支是正规子群，商群 $Diff(X)/Diff_0(X)$ 常被称为大微分同胚群。\footnote{更确切地说，大微分同胚是 $Diff(X)$ 中落到商群的非平凡元素上去的那些元素。} 在环面的情形，利用与万有覆盖之间的简单关系，可以证明 $Diff^+(T^2)/Diff^+_0(T^2)\cong SL(2,\mathbb{Z})$。
  % [Tutor Note: 「显然，这是环面的一个微分同胚」与「可以证明 $Diff^+(T^2)/Diff^+_0(T^2)\cong SL(2,\mathbb{Z})$」两处都被作者一句带过。要把 (2.141) 与 (2.139)(2.140) 对上，需注意 $\tilde\sigma=\delta\sigma^1-\gamma\sigma^2$ 这类线性变换保持 $\mathbb{Z}\oplus\mathbb{Z}$ 的格，且 $A$ 与 (2.141) 中的矩阵互为转置逆的重新排列；而商群的同构则是最简单情形下的映射类群（mapping class group）计算：$\pi_0 Diff^+(T^2)=\pi_0 Diff(T^2)=SL(2,\mathbb{Z})$，可用万有覆盖 $\mathbb{R}^2$ 上保持格的作用（Nielsen 实现）来验证。作者未给证明，属省略。]

  {\small\itshape ［原文标记：$\clubsuit$这条评注现在与 2.2.12 节的材料重复了。$\clubsuit$］}
  % [原文待办标记，原文 p.40：页边批注，位于评注 3 旁]

  \item 我们稍后会讨论这些正则变换／微分同胚如何作用在理论的 Hilbert 空间上。
\end{enumerate}

眼下，让我们取 $a_2$ 为坐标，于是测度成为

\begin{equation}
  \exp\left(\frac{i\kappa}{\pi}\int dt\,a_1\dot a_2\right)
  \tag{2.142}
  \label{eq:2.142}
\end{equation}

够简单的：量子化之后我们必须有

\begin{equation}
  [\hat a_1,\hat a_2]=-i\frac{\pi}{\kappa}
  \tag{2.143}
  \label{eq:2.143}
\end{equation}

我们用可归一化的波函数 $\psi(a_2)$ 描述量子态，其动量为

\begin{equation}
  p_{a_2}=\frac{\kappa}{\pi}a_1
  \tag{2.144}
  \label{eq:2.144}
\end{equation}

但现在我们得施加周期性条件：首先，$\psi(a_2)$ 在 $a_2\sim a_2+2\pi n$、$n\in\mathbb{Z}$ 之下必须是周期的。

然而，与粒子在圆上运动的量子化非常不同的是，我们还必须有 $a_2$ 的周期性。
% [Tutor Note: 原文此处第二个「周期性」应为 $a_1$：紧接的下一段就改口说「若位置空间有周期性，则动量被量子化」，并明确以 $a_1\sim a_1+2\pi$ 为条件，推出 (2.146) 的 $\exp(-i2\kappa\hat a_2)\psi=\psi$（因为按 (2.145)，$a_1\to a_1+\alpha$ 由 $\exp(-i\frac{\kappa}{\pi}\alpha\hat a_2)$ 生成，取 $\alpha=2\pi$ 即得该式）。这是原文的笔误，译文照原文保留，读者按 $a_1$ 理解即可。]

通常，位置空间的周期性意味着动量被量子化，从而 de Broglie 波才能装进这个周期性的空间。因此我们预期，动量的周期性应当把位置量子化。

确实，如果我们改用 $a_1$ 作坐标来量子化，那么

\begin{equation}
  \hat a_2=i\frac{\pi}{\kappa}\frac{\partial}{\partial a_1}
  \tag{2.145}
  \label{eq:2.145}
\end{equation}

于是 $\exp(-i\frac{\kappa}{\pi}\alpha\hat a_2)$ 把 $a_1$ 平移为 $a_1+\alpha$；因此，若我们希望有周期性 $a_1\sim a_1+2\pi$，那么回到以 $a_2$ 为坐标、以波函数 $\psi(a_2)$ 描述态的 Schrödinger 绘景，我们必须有

\begin{equation}
  e^{-i2\kappa\hat a_2}\psi=\psi
  \tag{2.146}
  \label{eq:2.146}
\end{equation}

但是，既然我们处在以 $a_2$ 为坐标的 Schrödinger 表示中，这就意味着

\begin{equation}
  e^{-i2\kappa a_2}\psi(a_2)=\psi(a_2)
  \tag{2.147}
  \label{eq:2.147}
\end{equation}

这意味着 $\psi$ 只能支撑在由 $a_2=\frac{n\pi}{\kappa}\bmod 2\pi\mathbb{Z}$ 给出的那些 $a_2$ 值上，于是不同的可能性为 $n=0,\dots,2\kappa-1 \bmod 2\kappa$。
% [Tutor Note: 相位与整性条件是本节的两个陷阱。(i) 相位：(2.145) 的 $i$、$\exp(-i\frac{\kappa}{\pi}\alpha\hat a_2)$ 中的 $-i$、以及 (2.146) 的 $e^{-i2\kappa\hat a_2}$，三者的符号必须配套。若把 (2.145) 写成 $-i\frac{\pi}{\kappa}\partial_{a_1}$，则 $\exp(-i\frac{\kappa}{\pi}\alpha\hat a_2)$ 平移的将是 $a_1\to a_1-\alpha$，此时要得到 $a_1\to a_1+2\pi$ 就得改用 $e^{+i2\kappa\hat a_2}$（最终格点集仍是 $a_2=n\pi/\kappa$）；指数与生成元的符号若不配套，中间步骤会自相矛盾。(ii) 整性：由 (2.147) 解出 $a_2=n\pi/\kappa$ 时，要求 $n$ 与 $n+2\kappa$ 给出同一点，故 $2\kappa$ 必须是整数；这正是预量子化线丛存在的整性条件（Weyl 整性），也是 (2.148) 求和上限为 $2\kappa-1$ 的原因。$\kappa$ 半整数时态空间维数为奇数，仍自洽；若 $\kappa$ 非半整数，则整套构造不成立。]

结论：既然动量与坐标都是周期的，一般的波函数形如

\begin{equation}
  \psi=\sum_{n=0}^{2\kappa-1}\psi_n\,\delta\left(a_2-\frac{n\pi}{\kappa}\right)
  \qquad \psi_r\in\mathbb{C}
  \tag{2.148}
  \label{eq:2.148}
\end{equation}
% [Tutor Note: 两点。(i) 原文右端的 $\psi_r\in\mathbb{C}$ 应为 $\psi_n\in\mathbb{C}$（求和指标是 $n$），此处照原文转写，未改。(ii) 作者跳步：这样写出的 $\psi$ 是 $\delta$ 函数的叠加，按 $L^2$ 意义并不可归一（原文自身在页边用「DOESN'T LOOK VERY NORMALIZABLE! JUSTIFY」点出了这一点）。正确的处理是几何量子化：预量子化线丛的平方可积全纯截面空间自然是 $2\kappa$ 维的，$\delta$ 基只是「位置表象」中的形式写法；相干态基即 (2.177) 中的 $\chi_{\mu,\kappa}$。]

{\small\itshape ［原文标记：$\clubsuit$看上去并不很可归一！请说明理由。$\clubsuit$］}
% [原文待办标记，原文 p.41：(2.148) 右侧页边批注]

我们于是得出结论：环面上的态空间 $\mathcal{H}(T^2)$ 是一个维数为 $2\kappa$ 的复 Hilbert 空间。

这个问题同构于「在环的离散近似上运动的粒子」的量子化：我们把环看作 $U(1)$，而「离散近似」则是循环群 $\mathbb{Z}_{2\kappa}$ 的嵌入。

\subsubsection{环面上平坦联络的量子化：相干态版本} % 原文: 2.2.9 Quantization Of Flat Connections On The Torus: The Coherent Story

本节中我们给拓扑曲面 $\Sigma_2$ 赋予一个复结构 $J$。$(\Sigma_2,J)$ 合起来定义一张 Riemann 面。

我们仍先考虑 $\Sigma_2$ 是环面的情形，如上一节。环面上的一般复结构所给出的 Riemann 面可以写成椭圆曲线：

\begin{equation}
  (\Sigma_2,J)=E_\tau=\mathbb{C}/(\mathbb{Z}\oplus\tau\mathbb{Z})
  \tag{2.149}
  \label{eq:2.149}
\end{equation}

其中 $\tau$ 是复数，且必须有非零虚部，以使商是一个好的流形。复结构由下述说法确定：

\begin{equation}
  d\zeta=d\sigma^1+\tau d\sigma^2
  \tag{2.150}
  \label{eq:2.150}
\end{equation}

是 $(1,0)$ 形式。我们需要在 $\Sigma_2$ 上选一个定向，我们选 $d\sigma^1\wedge d\sigma^2$。

{\small\itshape ［原文标记：$\clubsuit$换个更好的字体：这里需要一个花体的 z。$\clubsuit$］}
% [原文待办标记，原文 p.41：页边批注，作者想把符号换成一个「花体的 z」，故下文由 $\zeta$ 改记为 $z$]

给定复结构后，存在一个典范定向 $id\zeta\wedge d\bar\zeta$。把它算出来就会看到，我们应当取 $\operatorname{Im}\tau>0$。

数学事实：$\tau,\tau'$ 定义等价的复结构，当且仅当

\begin{equation}
  \tau'=\frac{a\tau+b}{c\tau+d} \qquad
  \begin{pmatrix}a & b\\ c & d\end{pmatrix}\in SL(2,\mathbb{Z})
  \tag{2.151}
  \label{eq:2.151}
\end{equation}

也就是说，$E_\tau$ 与 $E_{\tau'}$ 作为复流形同构，当且仅当 \eqref{eq:2.151} 成立。于是我们得到复结构的模空间（「栈」）

\begin{equation}
  \mathcal{M}_{cplx}(\Sigma_2)\cong\mathcal{H}/SL(2,\mathbb{Z})
  \tag{2.152}
  \label{eq:2.152}
\end{equation}

{\small\itshape ［原文标记：****／Keyhole：请用两个 2 阶、3 阶轨形点来解释「stack」一词。／****］}
% [原文待办标记，原文 p.42：作者留下一段「Keyhole」待补的说明]

$\Sigma_2$ 上的复结构诱导平坦规范场空间上的复结构：

\begin{equation}
  \begin{aligned}
    A&=a_1d\sigma^1+a_2d\sigma^2\\
    &=a_\zeta d\zeta+a_{\bar\zeta}d\bar\zeta\\
    a_\zeta&=\frac{a_2-\bar\tau a_1}{\tau-\bar\tau}\\
    a_{\bar\zeta}&=-\frac{a_2-\tau a_1}{\tau-\bar\tau}
  \end{aligned}
  \tag{2.153}
  \label{eq:2.153}
\end{equation}

$a_\zeta$ 与 $a_{\bar\zeta}$ 充当平坦联络模空间上的（反）全纯坐标。

定义略微修改过的坐标其实会很有用：

\begin{equation}
  \begin{aligned}
    \tilde a_{\bar\zeta}&=(a_2-\tau a_1)/2\pi\\
    \tilde a_\zeta&=(a_2-\bar\tau a_1)/2\pi
  \end{aligned}
  \tag{2.154}
  \label{eq:2.154}
\end{equation}

因为我们希望波函数随复结构全纯地变化，而不是反全纯地变化，\footnote{更概念化的理由是：Cauchy-Riemann 算子 $\bar\partial_A$ 在复结构变化下的改变为
\begin{equation}
  \bar\partial_A\to\bar\partial_A+\mu\partial_A
  \tag{2.155}
  \label{eq:2.155}
\end{equation}
其中 $\mu$ 是 Beltrami 微分：一个 $(-1,1)$ 型的形式。闭 Beltrami 微分构成的空间可以典范地等同于 $\Sigma_2$ 上复结构空间的全纯切空间。} 我们将取 $z=\tilde a_{\bar\zeta}$。我们把它归一化，使之具有周期性

\begin{equation}
  z\sim z+n+m\tau \qquad n,m\in\mathbb{Z}
  \tag{2.156}
  \label{eq:2.156}
\end{equation}

此时 Poisson 括号为

\begin{equation}
  \{z,\bar z\}=\frac{\operatorname{Im}\tau}{2\pi i}\frac{1}{\kappa}
  \tag{2.157}
  \label{eq:2.157}
\end{equation}

并且由 \eqref{eq:2.79}，我们有
% 跨片段引用说明：\eqref{eq:2.79} 指向 c03 片段中的 (2.79)（辛形式在模空间上的表达），本片段未重复引入；单独编译本片段时该引用会显示为「??」，属预期行为。

\begin{equation}
  \begin{aligned}
    \frac{\omega}{2\pi}
    &=\frac{\kappa}{(2\pi)^2}\int_{T^2}\delta A_s\wedge\delta A_s\\
    &=\frac{\kappa}{2\pi^2}\,\delta a_1\wedge\delta a_2\\
    &=i\kappa\,\frac{dz\wedge d\bar z}{\operatorname{Im}\tau}\\
    &=(2\kappa)\frac{i}{2\pi}\partial\bar\partial K
  \end{aligned}
  \tag{2.158}
  \label{eq:2.158}
\end{equation}

其中我们定义：

\begin{equation}
  K=2\pi\frac{(\operatorname{Im}z)^2}{\operatorname{Im}\tau}
  \tag{2.159}
  \label{eq:2.159}
\end{equation}

从而

\begin{equation}
  \partial\bar\partial K=\frac{\pi}{\operatorname{Im}\tau}dz\wedge d\bar z
  \tag{2.160}
  \label{eq:2.160}
\end{equation}
% [Tutor Note: 这里有两个 $2\pi$ 需要核对。其一，(2.158) 第二行的 $\kappa/2\pi^2$ 不是笔误：它等于 (2.137) 的 $\kappa/\pi$ 再除以 $2\pi$（左边是 $\omega/2\pi$）。其二，(2.158) 末行与 (2.160) 的相容性依赖于约定 $\partial\bar\partial K\equiv(\partial_z\partial_{\bar z}K)\,dz\wedge d\bar z$；用 (2.159) 直接算得 $\partial_z\partial_{\bar z}K=2\pi/\operatorname{Im}\tau\cdot\frac12=\pi/\operatorname{Im}\tau$（因为 $(\operatorname{Im}z)^2=-(z-\bar z)^2/4$），与 (2.160) 一致，再代入末行确实回到 $i\kappa\,dz\wedge d\bar z/\operatorname{Im}\tau$。若把 $\partial\bar\partial$ 的归一化取成 $\frac{i}{2\pi}\partial\bar\partial$（Kähler 形式的习惯），(2.158) 末尾的系数 $(2\kappa)$ 就会改变，务必逐项核对。]

导出波函数的方法有很多，其中之一如下：Berezin-Toeplitz 量子化告诉我们，去找一个全纯线丛 $\mathcal{L}_\kappa\to E_\tau$，配备一个不变 Hermitian 度规

\begin{equation}
  \|s(z)\|^2=e^{-K_\kappa}|s(z)|^2
  \tag{2.161}
  \label{eq:2.161}
\end{equation}

其中在差一个 Kähler 变换的意义下，

\begin{equation}
  K_\kappa=(2\kappa)K=4\pi\kappa\frac{(\operatorname{Im}z)^2}{\operatorname{Im}\tau}
  \tag{2.162}
  \label{eq:2.162}
\end{equation}

Hilbert 空间上的内积为

\begin{equation}
  \langle\psi_1,\psi_2\rangle
  =\int_{E_\tau}e^{-K_\kappa}\psi_1^*(z)\psi_2(z)\,\frac{\omega}{2\pi}
  \tag{2.163}
  \label{eq:2.163}
\end{equation}

其中我们用了线丛的局部平凡化，从而把 $\psi_1(z)$ 与 $\psi_2(z)$ 看作局部定义的全纯函数。等价地：我们把它拉回到万有覆盖，并在 $\mathbb{C}$ 中对 $\mathbb{Z}+\tau\mathbb{Z}$ 的平移作用取一个基本域作积分。

{\small\itshape ［原文标记：$\clubsuit$内积的整体归一化因子？？？？$\clubsuit$］}
% [原文待办标记，原文 p.43：页边批注，作者对内积的整体归一化未作确定]

考虑一个基本的全纯线丛 $\mathcal{L}_\theta\to E_\tau$ 会很有用，其中我们采用度规为

于是我们可以取 $\mathcal{L}_\kappa=\mathcal{L}_\theta^{\otimes 2\kappa}$。

那么，该如何描述 $\mathcal{L}_\theta$ 的全纯截面呢？嗯，

\begin{equation}
  e^{-2\pi\frac{(\operatorname{Im}z)^2}{\operatorname{Im}\tau}}|\psi(z)|^2
  \tag{2.164}
  \label{eq:2.164}
\end{equation}

必须是良定义的。
% [排版说明：原文把 (2.164) 的度规表达式错放在「Well,」之后（上一句「where we use the metric with」被截断）。这正是一个度规：$e^{-2\pi(\operatorname{Im}z)^2/\operatorname{Im}\tau}|\psi(z)|^2$，它等于 (2.162) 在 $2\kappa=1$（即 $\kappa=1/2$，此时 $\mathcal{L}_\kappa=\mathcal{L}_\theta$）时的情形，与下文「$\mathcal{L}_\kappa=\mathcal{L}_\theta^{\otimes 2\kappa}$」相容。译文按原文语序保留该位置。]

于是我们可以把 $\psi(z)$ 提升为万有覆盖上（即复 $z$ 平面上）的整函数，然后要求在 $z\to z+a+b\tau$（$a,b\in\mathbb{Z}$）之下不变。现在，在 $z\to z+a$ 之下，$\psi(z)$ 至多只能改变一个相位，而由于它是全纯的，这个相位必须为 1。另一方面，在 $z\to z+b\tau$ 之下，一个简短的计算表明

\begin{equation}
  \begin{aligned}
    e^{-2\pi\frac{(\operatorname{Im}(z+b\tau))^2}{\operatorname{Im}\tau}}
    &=e^{-2\pi\frac{(\operatorname{Im}z)^2}{\operatorname{Im}\tau}}
    e^{2\pi ib(z-\bar z)+i\pi b^2(\tau-\bar\tau)}\\
    &:=e^{-2\pi\frac{(\operatorname{Im}z)^2}{\operatorname{Im}\tau}}
    \,e_b(z,\tau)^{-1}\,\overline{e_b(z,\tau)}^{-1}
  \end{aligned}
  \tag{2.165}
  \label{eq:2.165}
\end{equation}
% [Tutor Note: (2.165) 第二行第二个因子上有一道横线（共轭），PDF 文本层把横线丢掉了，只留下 $e_b^{-1}e_b^{-1}$ 的样子。按 (2.166) 的 $e_b(z,\tau)=e^{-i\pi b^2\tau-2\pi ibz}$，有 $e_b^{-1}=e^{i\pi b^2\tau+2\pi ibz}$、$\overline{e_b}^{-1}=e^{-i\pi b^2\bar\tau-2\pi ib\bar z}$，两者相乘恰为 $e^{2\pi ib(z-\bar z)+i\pi b^2(\tau-\bar\tau)}$；漏掉共轭就无法还原第一行。此处已按页面图像补回横线。]

其中

\begin{equation}
  e_b(z,\tau)=e^{-i\pi b^2\tau-2\pi ibz}
  \tag{2.166}
  \label{eq:2.166}
\end{equation}

再次利用全纯性排除可能多出的相位依赖，我们得知：截面到万有覆盖的提升必须是一个整函数，它是准周期的，并满足函数方程

\begin{equation}
  \psi(z+a+b\tau)=e_b(z,\tau)\psi(z) \qquad a,b\in\mathbb{Z}
  \tag{2.167}
  \label{eq:2.167}
\end{equation}

我们将记 $e_1(z,\tau)=e(z,\tau)$。这个函数方程很容易解：在 $z\to z+1$ 下的不变性意味着 $\psi(z)$ 必须有 Fourier 展开

\begin{equation}
  \psi(z)=\sum_{n\in\mathbb{Z}}\psi_n e^{2\pi inz}
  \tag{2.168}
  \label{eq:2.168}
\end{equation}

而在 $z\to z+\tau$ 下的变换则给出

\begin{equation}
  \psi_{n+1}=e^{i\pi(2n+1)\tau}\psi_n \qquad n\in\mathbb{Z}
  \tag{2.169}
  \label{eq:2.169}
\end{equation}

从而 $\psi_n=\psi_0e^{i\pi n^2\tau}$。

因此，在这个基本情形中，截面空间是一维的，它由 Riemann theta 函数张成：

\begin{equation}
  \theta(z,\tau)=\sum_{n\in\mathbb{Z}}e^{i\pi n^2\tau+2\pi inz}
  \tag{2.170}
  \label{eq:2.170}
\end{equation}

虽然我们把 $\theta(z,\tau)$ 写成了 $z$ 平面上的整函数，它也可以看作某个全纯线丛的全纯截面。

{\small\itshape ［原文标记：$\clubsuit$重复了吗？$\clubsuit$］}
% [原文待办标记，原文 p.44：页边批注]

接着我们定义与 Chern 类为 1 的主 $U(1)$ 丛相伴的复线丛

\begin{equation}
  L=(\mathbb{C}\times V)/(\mathbb{Z}\oplus\mathbb{Z})
  \tag{2.171}
  \label{eq:2.171}
\end{equation}

其中我们把 $V$ 看作 $U(1)$ 的一维表示。$\mathbb{Z}\oplus\mathbb{Z}$ 的一对生成元 $(g_1,g_2)$ 作用为

\begin{equation}
  \begin{aligned}
    g_1&:(z,\psi)\to(z+1,\psi)\\
    g_2&:(z,\psi)\to(z+\tau,e(z,\tau)\psi)
  \end{aligned}
  \tag{2.172}
  \label{eq:2.172}
\end{equation}

若取 $V_\ell$ 为 $V^{\otimes\ell}$，即 $U(1)$ 的荷 $\ell$ 表示，则我们得到

\begin{equation}
  L^{\otimes\ell}=(\mathbb{C}\times V_\ell)/(\mathbb{Z}\oplus\mathbb{Z})
  \tag{2.173}
  \label{eq:2.173}
\end{equation}

此时生成元作用为

\begin{equation}
  \begin{aligned}
    g_1&:(z,\psi)\to(z+1,\psi)\\
    g_2&:(z,\psi)\to(z+\tau,e(z,\tau)^\ell\psi)
  \end{aligned}
  \tag{2.174}
  \label{eq:2.174}
\end{equation}

因此，$L^{\otimes\ell}$ 的全纯截面满足如下恒等式：

\begin{equation}
  \begin{aligned}
    \psi(z+1)&=\psi(z)\\
    \psi(z+\tau)&=e^{-\ell\pi i\tau-2\ell\pi iz}\psi(z)
  \end{aligned}
  \tag{2.175}
  \label{eq:2.175}
\end{equation}

当 $\ell\ge 0$ 时存在整体全纯截面。当 $\ell<0$ 时没有整体全纯截面。

现在，正如我们所说，在量子化能级为 $\kappa$ 的 $U(1)$ Chern-Simons 理论时，相关的全纯线丛是 $\mathcal{L}_\theta^{2\kappa}$，因此若 $\psi(z)$ 是波函数，则

\begin{equation}
  \langle\psi_1,\psi_2\rangle
  =\int_{E_\tau}e^{-4\pi ik\frac{(\operatorname{Im}z)^2}{\operatorname{Im}\tau}}
  \psi_1^*(z)\psi_2(z)\frac{\omega}{2\pi}
  \tag{2.176}
  \label{eq:2.176}
\end{equation}
% [Tutor Note: (2.176) 的指数印作 $e^{-4\pi i k(\operatorname{Im}z)^2/\operatorname{Im}\tau}$，其中既有虚数单位 $i$，又把 $\kappa$ 写成了 $k$。对照 (2.162) 的 $K_\kappa=4\pi\kappa(\operatorname{Im}z)^2/\operatorname{Im}\tau$（实）与 (2.163) 的 $e^{-K_\kappa}$，此处应为实指数 $e^{-4\pi\kappa(\operatorname{Im}z)^2/\operatorname{Im}\tau}$：多出的 $i$ 会使被积函数变成纯相位，$2\kappa$ 维截面空间上的内积就不再是 Hermitian 的。原文疑有笔误，此处照原文转写，未改。]

该线丛有 $c_1(\mathcal{L}_\theta^{2\kappa})=2\kappa\,c_1(\mathcal{L}_\theta)=2\kappa$，因而将存在 $2\kappa$ 维的全纯截面空间。一组自然的截面基由能级 $\kappa$ 的 theta 函数 $\Theta_{\mu,\kappa}$ 给出。关于能级 $\kappa$ 的 theta 函数的更多细节见附录 D。

\textbf{评注}：这里我们假定 $\kappa>0$。若翻转 $\Sigma_2$ 的定向，那就会改变 $\mathcal{M}_{flat}$ 的定向，此时需要取 $\kappa<0$ 才有非零截面。

总之，当我们额外选定 $\Sigma_2$ 上的一个复结构 $J$ 作为附加数据时，我们就能构造 Chern-Simons 理论的物理态空间 $\mathcal{H}(\Sigma_2,J)$。在全纯量子化中，环面上物理态的一组基是

\begin{equation}
  \chi_{\mu,\kappa}(z,\tau)=\frac{\Theta_{\mu,\kappa}(z,\tau)}{\eta(\tau)}
  \qquad \mu\sim\mu+2\kappa
  \tag{2.177}
  \label{eq:2.177}
\end{equation}

$1/\eta(\tau)$ 这个因子在目前这一步完全不显然。我们将在下面讨论它。它可以通过以下方式得到：

\begin{enumerate}
  \item 沿复结构模空间作平行移动。
  \item 在「所有规范场构成的空间」的量子化中（并受 Gauss 定律约束）对波泛函 $\Psi[A]$ 作归一化。
\end{enumerate}

\textbf{重要评注}：用这种量子化方式，我们得到的态空间 $\mathcal{H}(\Sigma_2,J)$ 依赖于 $\Sigma_2$ 上的复结构 $J$。但 $J$ 并不是问题数据的一部分！在拓扑场论中，我们必须「摆脱」对 $J$ 的依赖。我们应把 $\mathcal{H}(\Sigma_2,J)$ 看作 $\mathcal{M}_{cplx}(\Sigma_2)$ 上某个向量丛的纤维。「态空间与 $J$ 无关」这一陈述，是由该丛上存在一个射影平坦联络来实现的。关于那个联络，详见下文。

{\small\itshape ［原文标记：****／应当讨论协变导数：上述 Hermitian 度规的 Chern 联络。／****］}
% [原文待办标记，原文 p.45：作者留待补写的段落]

\subsubsection*{推广到亏格 $g$ 的紧 Riemann 面} % 原文: Generalization To Compact Riemann Surface Of Genus g（原文为无编号、带下划线的小标题，故用 \subsubsection* 复刻）

相干态方法的一个优点是，它很容易推广到亏格为 $g$ 的紧定向曲面 $\Sigma_2$。

回忆相空间可以等同于 $2g$ 维环面 $\mathcal{P}=H^1(\Sigma;\mathbb{R})/H^1(\Sigma;\mathbb{Z})$。

若给 $\Sigma_2$ 一个复结构，使它成为 Riemann 面，那么正如我们上面显式描述过的环面情形，$\mathcal{M}_{flat}$ 继承一个复结构，满足

\begin{equation}
  T^{1,0}\mathcal{M}_{flat}\cong H^{0,1}(\Sigma_2)
  \tag{2.178}
  \label{eq:2.178}
\end{equation}

于是，我们的环面 $\mathcal{P}\cong(S^1)^{2g}$ 现在是一个复环面。预量子线丛 $L_k\to\mathcal{M}_{flat}$ 成为 $L_k=\mathcal{L}_\theta^{\otimes 2k}$，其中 $\mathcal{L}_\theta$ 仍是标准全纯线丛，具有一维的全纯截面空间。这属于一个一般的故事：（更多细节见附录 E。）

\begin{enumerate}
  \item 给定带整辛形式的辛格（symplectic lattice），可以典范地构造一个 Heisenberg 群，作为 $\Lambda\otimes\mathbb{R}/\Lambda$ 的 $U(1)$ 扩张。
  \item 该 Heisenberg 群是环面 $\Lambda\otimes\mathbb{R}/\Lambda$ 上的主 $U(1)$ 丛。
  \item 要表示该 Heisenberg 群，我们需要选取格的一个拉格朗日分解。
  \item 此外，若我们还有一个相容的复结构，那么由 $U(1)$ 的定义表示与 Heisenberg 群相伴的复线丛就是一个带典范联络的全纯线丛，其曲率代表 $c_1(L)$。
  \item 全纯截面可以提升到万有覆盖 $\Lambda\otimes\mathbb{R}\cong\mathbb{C}^n$，并且就复结构而言它是一个整函数。它是准周期函数，其自守因子由 Heisenberg 群的余循环决定。这就定义了基本的 theta 函数。
  \item $\mathcal{L}_\theta^{\otimes 2\kappa}$ 的全纯截面的一组基，是相对于 $\Lambda$ 的一个拉格朗日分裂而定义的能级 $\kappa$ 的 theta 函数。
\end{enumerate}

在我们的情形，辛格就是 $H_1(\Sigma;\mathbb{Z})$，辛形式由相交形式提供。（请记住：$\Sigma$ 是定向的。）于是 $\Sigma$ 上的复结构如前所述诱导 $H_1(\Sigma;\mathbb{R})$ 上的复结构。剩下的只是把机器开动起来。

非常具体地说：在 $\Sigma$ 上取 A、B 圈的一组基，以及对偶 1-形式 $\alpha_I,\beta^I$，$I=1,\dots,g$。于是存在一组调和 $(0,1)$ 形式：

\begin{equation}
  f_I:=\alpha_I+\tau_{IJ}\beta^J \qquad I=1,\dots,g
  \tag{2.179}
  \label{eq:2.179}
\end{equation}

而调和 $(1,0)$ 形式由下式张成

\begin{equation}
  \bar f_I:=\alpha_I+\bar\tau_{IJ}\beta^J \qquad I=1,\dots,g
  \tag{2.180}
  \label{eq:2.180}
\end{equation}
% [约定说明：原文把 $f_I=\alpha_I+\tau_{IJ}\beta^J$ 称作 $(0,1)$ 形式、把 $\bar f_I$ 称作 $(1,0)$ 形式，这与 (2.150) 处「$d\zeta=d\sigma^1+\tau d\sigma^2$ 是 $(1,0)$ 形式」的约定正好相反；按 (2.150) 的约定以及周期矩阵 $\tau_{IJ}$ 的通常用法，$f_I$ 应是全纯 $(1,0)$ 形式（附录 E 的 theta 函数即以此为基础）。此处照原文转写，读者按 (2.150) 的约定理解即可，勿与 2.2.9 节环面情形的 $\zeta$／$z$ 记号混用。]
